\documentclass[11pt,reqno]{amsart}

\usepackage[margin=1.08in]{geometry}
\usepackage{amsmath,amssymb,amsthm,mathtools}
\usepackage{booktabs}
\usepackage{tabularx}
\usepackage{microtype}
\usepackage{enumitem}
\usepackage{xcolor}
\usepackage[hidelinks]{hyperref}
\usepackage[nameinlink,capitalise,noabbrev]{cleveref}

\hypersetup{
  pdftitle={A Sharper Numerical Certificate for the Matrix Multiplication Exponent},
  pdfauthor={Alex Townsend}
}

\newtheorem{theorem}{Theorem}[section]
\newtheorem{proposition}[theorem]{Proposition}
\newtheorem{lemma}[theorem]{Lemma}
\newtheorem{corollary}[theorem]{Corollary}
\theoremstyle{definition}
\newtheorem{definition}[theorem]{Definition}
\theoremstyle{remark}
\newtheorem{remark}[theorem]{Remark}

\newcommand{\CW}{\mathrm{CW}}
\newcommand{\MM}{\langle\!\langle}
\newcommand{\MMend}{\rangle\!\rangle}
\newcommand{\Hent}{\mathrm{H}}
\newcommand{\Rbar}{\widetilde{R}}
\newcommand{\eps}{\varepsilon}
\newcommand{\DeltaSix}{\Delta_{6}}

\title[A sharper certificate for $\omega$]{A Sharper Numerical Certificate for the\\ Matrix Multiplication Exponent}

\author{Alex Townsend}
\address{Department of Mathematics, Cornell University, Ithaca, NY 14853, USA}
\email{townsend@cornell.edu}
\date{August 9, 2026}

\subjclass[2020]{68Q17, 15A69, 68W40}
\keywords{matrix multiplication exponent, Coppersmith--Winograd tensor, laser method, computer-assisted proof, constrained optimisation}

\begin{document}

\begin{abstract}
We report a new computer-assisted world-record candidate for the exponent of square matrix multiplication,
\[
  \omega < 2.371310.
\]
The previous best published bound was $\omega<2.371339$, obtained from the fourth power of the
Coppersmith--Winograd tensor by the \emph{More Asymmetry} analysis of Alman, Duan,
Vassilevska Williams, Xu, Xu, and Zhou.  We do not change their combinatorial degeneration.
Instead, we re-optimise 126 term-specific six-region probability vectors that are already independent
variables in their constituent-stage theorem.  Replacing every nonsmooth minimum by an exact epigraph
gives 42 explicit candidate inequalities and exposes a better balanced feasible point.  A fresh evaluation
gives the tightened exponent $2.3713098963342496$; a conservative maximum-entropy dual audit gives
$2.3713098963360215$.  We supply a standalone certificate with stored exponent
$2.3713099046652784$, inward inequality slack, an independent NumPy verifier, and a second
term-by-term reconstruction.  Finally, tensoring the fourth-power degeneration with itself gives an exact
eighth-power realisation with exponent $2.3713098963357147$.  The complete certificate and audit trail
are included as ancillary files.  The analytic reduction is exact; the supplied validation is a
reproducible, dual-safe binary64 computation rather than an outward-rounded interval proof.
\end{abstract}

\maketitle

\section{Introduction}

Let $\omega$ denote the infimum of the real numbers $\tau$ such that two $n\times n$ matrices over a
field can be multiplied using $O(n^{\tau+\eps})$ arithmetic operations for every $\eps>0$.
Strassen's seven-multiplication algorithm showed in 1969 that $\omega<\log_2 7$~\cite{Strassen1969}.
The subsequent record sequence has been driven by increasingly refined analyses of powers of the
Coppersmith--Winograd tensor~\cite{CoppersmithWinograd1990}.  The most recent steps are shown in
\cref{tab:records}.

\begin{table}[t]
  \centering
  \caption{Recent improvements to the square matrix multiplication exponent.  The last row is the
  conservative decimal statement supported by the certificate in this paper.}
  \label{tab:records}
  \small
  \begin{tabularx}{\textwidth}{@{}>{\raggedright\arraybackslash}p{0.30\textwidth}Xr@{}}
    \toprule
    Work & Main mechanism & Bound on $\omega$ \\
    \midrule
    Alman--Vassilevska Williams~\cite{AlmanWilliams2024}
      & refined laser method & $2.3728596$ \\
    Duan--Wu--Zhou~\cite{DuanWuZhou2023}
      & asymmetric hashing, eighth power & $2.371866$ \\
    Vassilevska Williams et al.~\cite{VXXZ2024}
      & complete split distributions & $2.371552$ \\
    Alman et al.~\cite{MoreAsymmetry}
      & six-region More Asymmetry analysis & $2.371339$ \\
    This paper
      & exact epigraph rebalancing and certification & $2.371310$ \\
    \bottomrule
  \end{tabularx}
\end{table}

As of August 9, 2026, the best published bound, $\omega<2.371339$, follows from the \emph{More Asymmetry}
framework~\cite{MoreAsymmetry}.  That work analyses $\CW_5^{\otimes4}$ and introduces six possible
permutations of the roles of the $X$, $Y$, and $Z$ variables.  At every interior constituent term, the
theorem permits an independent probability vector describing how much of that term is assigned to each
of the six regions.  The released implementation registers these vectors independently, but the resulting
optimisation is large, nonsmooth, and numerically delicate.  The authors report that a typical solve takes
about a week and that the full eighth-power programme was infeasible with their solver.

Our contribution is deliberately narrow but consequential: we freeze the published global distributions,
conditional split laws, complete split distributions, and maximum-entropy witnesses, and re-optimise only
the 126 exposed constituent-stage region vectors.  This is enough to improve the exact stored record point
from $2.3713389005434182$ to the safe value $2.3713099046652784$.  The rounded improvement is
approximately $2.90\times10^{-5}$.

The main result is as follows.

\begin{theorem}[Record-candidate bound]\label{thm:main}
  The square matrix multiplication exponent satisfies
  \[
     \omega<2.371310.
  \]
\end{theorem}

Three features of the computation are worth separating.

\begin{enumerate}[leftmargin=2.2em]
  \item \emph{The structural theorem is unchanged.}  Every modified quantity is a six-entry simplex
  variable explicitly allowed by the constituent-stage theorem of~\cite{MoreAsymmetry}.
  \item \emph{The discovery solve is not the certificate.}  We use an exact epigraph formulation to find
  the row point, then discard its auxiliary variables and reconstruct every retained-copy and matrix-size
  quantity in a fresh model with explicit inward slack.
  \item \emph{The eighth-power statement is exact.}  Tensoring the certified fourth-power degeneration
  with itself gives a degeneration of $\CW_5^{\otimes8}$ with the same exponent.  No conjectural
  correlated eighth-power improvement is needed for \cref{thm:main}.
\end{enumerate}

The remainder of the paper gives enough of the framework to identify the modified variables, states the
finite optimisation problem, and records three independent certification paths.  Full parameter arrays are
left to the machine-readable ancillary files rather than printed in the paper.

\section{Tensor background and the More Asymmetry programme}

\subsection{Matrix multiplication tensors}

For positive integers $a,b,c$, write
\[
  \langle a,b,c\rangle
  =\sum_{i=1}^{a}\sum_{j=1}^{b}\sum_{k=1}^{c}x_{ij}y_{jk}z_{ki}
\]
for the matrix multiplication tensor.  We write $\Rbar(T)$ for the asymptotic rank of a tensor $T$.
Sch\"onhage's asymptotic sum inequality~\cite{Schonhage1981} converts a degeneration of a low-rank
tensor into a large direct sum of matrix multiplication tensors into an upper bound on $\omega$.

The Coppersmith--Winograd tensor is
\begin{align*}
  \CW_q={}&x_0y_0z_{q+1}+x_0y_{q+1}z_0+x_{q+1}y_0z_0 \\
  &+\sum_{i=1}^{q}\bigl(x_0y_iz_i+x_iy_0z_i+x_iy_iz_0\bigr).
\end{align*}
It has border rank and asymptotic rank $q+2$~\cite{CoppersmithWinograd1990}.  The recent record bounds
take $q=5$, recursively zero out powers of $\CW_5$ into direct sums of matrix multiplication tensors,
and apply the asymptotic sum inequality.

\subsection{The finite bound at the fourth power}

We use the notation of~\cite[Sections 5--7]{MoreAsymmetry}.  At recursion level
$\ell^*=3$, the input tensor is $\CW_5^{\otimes 2^{\ell^*-1}}=\CW_5^{\otimes4}$.
The global stage applies Theorem 5.3 of that paper, and the constituent stages apply its Theorems 6.2
and 6.4.  The result is a degeneration of many copies of $\CW_5^{\otimes4n}$ into a direct sum of
matrix multiplication tensors.  After taking logarithms, the sufficient condition in
\cite[Equation (7)]{MoreAsymmetry} has the form
\begin{equation}\label{eq:master}
   \mathcal R(\theta)+\tau\,\mathcal M(\theta)\ge 4\log 7.
\end{equation}
Here $\theta$ denotes all probability distributions and complete-split data, $\mathcal R$ is the total
logarithmic retained-copy contribution, and $\mathcal M$ is the logarithm of the side length of the final
square matrix multiplication tensor.  Any feasible parameter point satisfying \cref{eq:master} implies
$\omega\le\tau$.

For the square problem, the six global regions have equal mass $1/6$.  At every interior constituent term
$t$, however, Theorem 6.4 permits its own vector
\begin{equation}\label{eq:simplex}
  \rho_t=(\rho_{t,1},\ldots,\rho_{t,6})\in\DeltaSix,
  \qquad
  \rho_{t,r}\ge0,
  \qquad
  \sum_{r=1}^{6}\rho_{t,r}=1.
\end{equation}
The six entries select the six permutations of the roles of $X,Y,Z$.  The released fourth-power model
contains 126 interior terms with such vectors.  Their conditional split distributions remain fixed in this
paper; only the proportions in \cref{eq:simplex} vary.

For later use, expose the three coordinate candidates before each minimum:
\begin{itemize}[leftmargin=2em]
  \item $L^{(3)}_{r,d}(\rho)$, for the six level-three region retentions;
  \item $L^{(2)}_d(\rho)$, for the symmetric level-two retention;
  \item $G_{r,d}(\rho)$, for the six global-stage retentions; and
  \item $M_d(\rho)$, for the three output matrix-size candidates,
\end{itemize}
where $r\in[6]$ and $d\in\{X,Y,Z\}$.  These are exactly the unminimised arrays computed by the released
verification model.  The total retained and matrix-size quantities are
\begin{align}
 \mathcal R(\rho)
  &=\sum_{r=1}^{6}\min_d L^{(3)}_{r,d}(\rho)
    +\min_d L^{(2)}_d(\rho)
    +\sum_{r=1}^{6}\min_d G_{r,d}(\rho),\label{eq:retained}\\
 \mathcal M(\rho)&=\min_d M_d(\rho).\label{eq:matrix}
\end{align}
Thus the exponent certified by a structural row point is
\begin{equation}\label{eq:omega-ratio}
  \Omega(\rho)=\frac{4\log 7-\mathcal R(\rho)}{\mathcal M(\rho)}.
\end{equation}

\begin{lemma}[Legality of the optimisation axis]\label{lem:legal}
Fix every parameter in the published $\CW_5^{\otimes4}$ certificate except the 126 vectors
$\{\rho_t\}$.  If each $\rho_t$ satisfies \cref{eq:simplex}, and if all mixed complete split distributions,
compatibility terms, minima, and dependent auxiliary variables are recomputed, then the resulting point is
admissible in Theorems 5.3, 6.2, and 6.4 of~\cite{MoreAsymmetry}.
\end{lemma}

\begin{proof}
For every interior constituent term $t$, Theorem 6.4 assigns independent nonnegative quantities
$A_{t,r}$ with $\sum_r A_{t,r}=1$.  These are precisely the entries stored as the term's local region
proportions.  Changing $A_{t,r}$ repartitions occurrences of the already revealed parent term among the
six hashing orientations.  It does not change the parent identifier, a conditional split law, or a
maximum-entropy feasible set.  For fixed lower policies, the level-three, level-two, and matrix-size
contributions are affine in the region weights.  The positional complete split distribution is affine before
the global entropy and compatibility functionals are evaluated.  Recomputing those nonlinear functionals
and all subsequent minima therefore evaluates the same theorem at a different admissible point.
\end{proof}

This lemma is the conceptual reason the improvement does not require a new laser-method argument.
The work is in locating and carefully certifying a better point of the existing finite programme.

\section{Exact epigraph rebalancing}

\subsection{Removing the nonsmooth minima}

Direct smooth-min approximations proved useful for exploration but obscured which coordinate branches
were active.  We therefore introduced 14 epigraph variables
\[
 e^{(3)}_1,\ldots,e^{(3)}_6,\quad e^{(2)},\quad
 g_1,\ldots,g_6,\quad m
\]
and imposed the 42 inequalities
\begin{align}
 e^{(3)}_r&\le L^{(3)}_{r,d}(\rho), &&r\in[6],\ d\in\{X,Y,Z\},\label{eq:epi3}\\
 e^{(2)}&\le L^{(2)}_d(\rho), &&d\in\{X,Y,Z\},\label{eq:epi2}\\
 g_r&\le G_{r,d}(\rho), &&r\in[6],\ d\in\{X,Y,Z\},\label{eq:epig}\\
 m&\le M_d(\rho), &&d\in\{X,Y,Z\}.\label{eq:epim}
\end{align}
At any optimum the appropriate inequalities are tight, and \cref{eq:omega-ratio} becomes
\begin{equation}\label{eq:epi-objective}
  \min_{\rho,e,g,m}
  \frac{4\log 7-\sum_{r=1}^{6}e^{(3)}_r-e^{(2)}-\sum_{r=1}^{6}g_r}{m}
\end{equation}
subject to \cref{eq:simplex,eq:epi3,eq:epi2,eq:epig,eq:epim}.  The formulation has 756 region entries
and 14 epigraph variables.  We parameterised each supported simplex by a masked softmax and supplied
analytic derivatives of every objective and constraint.

The optimisation remains nonconvex: in particular, the global complete-split entropy depends nonlinearly
on the mixed region rows.  Global optimality is neither claimed nor needed.  Any point returned by the
search is evaluated independently, and a single feasible improvement suffices.

\subsection{Search and reconstruction}

Starting from the published certificate, we first used smooth continuation to approach the active-set
intersection.  We then solved \cref{eq:epi-objective} by sequential quadratic programming, using the SLSQP
implementation in SciPy~\cite{Kraft1988,SciPy2020}.  The final scaled active-set stage converged in two iterations;
the minimum epigraph residual was zero at displayed precision, and the epigraph and direct objectives
agreed to $3\times10^{-15}$.

The solver's auxiliary variables are not trusted as a certificate.  We retain only the 126 row vectors,
install them into a fresh copy of the released model, recompute every complete split distribution and
unminimised candidate, take the exact coordinate minima, and rebuild the retained-copy, matrix-size, and
exponent variables with inward slack.  The numerical progression is summarised in
\cref{tab:values}.

\begin{table}[t]
  \centering
  \caption{Exponent values associated with the final structural row point.}
  \label{tab:values}
  \begin{tabular}{@{}lr@{}}
    \toprule
    Quantity & Value \\
    \midrule
    Previous stored certificate & $2.3713389005434182$ \\
    Direct epigraph-row evaluation & $2.371309896335721$ \\
    Fresh verifier, exact coordinate minima & $2.3713098963342496$ \\
    Maximum-entropy dual-safe evaluation & $2.3713098963360215$ \\
    Stored certificate with safety slack & $2.3713099046652784$ \\
    Conservative public statement & $2.371310$ \\
    \bottomrule
  \end{tabular}
\end{table}

The change audit finds 543 modified scalar entries among the 24,855 parameters: 528 coordinates in the
126 legal region rows and 15 reconstructed auxiliary coordinates.  No global distribution, split law,
complete split distribution, or entropy witness is changed.

\section{Certification}

\subsection{A dual-safe entropy calculation}

The More Asymmetry programme contains many penalties
\[
  P_\alpha=\max\{\Hent(p):Ap=\mu,\ p\ge0\}-\Hent(\alpha),
  \qquad
  \Hent(p)=-\sum_s p_s\log p_s.
\]
The released certificate stores both an approximate maximum-entropy primal distribution and KKT
multipliers.  We remove any logical dependence on the approximate KKT equalities by using weak duality.

\begin{lemma}[Entropy dual bound]\label{lem:entropy-dual}
For every vector $\lambda$ and every feasible marginal vector $\mu$,
\begin{equation}\label{eq:entropy-dual}
 \max_{p\ge0,\,Ap=\mu}\Hent(p)
 \le D(\lambda;\mu)
 :=\sum_s \exp\bigl((A^\mathsf{T}\lambda)_s-1\bigr)-\lambda^\mathsf{T}\mu.
\end{equation}
\end{lemma}

\begin{proof}
Use the Lagrangian
$\Hent(p)+\lambda^\mathsf{T}(Ap-\mu)$ and the scalar identity
\[
  \sup_{x\ge0}\{-x\log x+yx\}=e^{y-1}.
\]
Taking the supremum over all entries of $p$ gives \cref{eq:entropy-dual}; weak duality completes the
proof.
\end{proof}

For every one of the 762 maximum-entropy subproblems, we evaluate $D$ at the stored multipliers and
replace the primal entropy by this upper bound.  This can only increase $P_\alpha$, decrease the relevant
retention candidate, and hence increase the inferred exponent.  The resulting conservative value is
\[
  \omega_{\rm dual}=2.3713098963360215,
\]
which remains $1.0366\times10^{-7}$ below the stated decimal bound.  The approximate KKT residuals are
therefore diagnostics, not proof assumptions.  The dual-safe logarithmic rates are
\[
  \mathcal M=2.0967467985574295,
  \qquad
  \mathcal R_{\rm dual}=2.8116041626911494,
\]
so that $\omega_{\rm dual}=(4\log 7-\mathcal R_{\rm dual})/\mathcal M$.

\subsection{Independent evaluations}

We performed three complementary checks.

\begin{enumerate}[leftmargin=2.2em]
  \item A fresh value-only NumPy port reconstructs the 24,855-parameter model from the MAT file.  It
  does not use SNOPT or the released forward automatic-differentiation class.
  \item A second dense checker explicitly drives every term's initialisation, downward mass propagation,
  and upward contribution propagation, then independently sums every unminimised level-three,
  level-two, global, and matrix candidate.  It does not call the aggregate objective routine.
  \item A complete parameter-difference audit compares the new MAT file with the published source
  certificate group by group and confirms that every structural change lies in one of the 126 permitted
  region simplexes.
\end{enumerate}

\begin{table}[t]
  \centering
  \caption{Numerical certification summary.  Inequality residuals use the convention that positive values
  violate feasibility.}
  \label{tab:audit}
  \begin{tabular}{@{}lr@{}}
    \toprule
    Check & Result \\
    \midrule
    Parameters reconstructed & $24{,}855$ \\
    Maximum inequality residual & $-9.9999997\times10^{-10}$ \\
    Sch\"onhage inequality residual & $-2.0967486\times10^{-9}$ \\
    Maximum bound violation & $0$ \\
    Maximum optimised-row simplex residual & $2.22\times10^{-16}$ \\
    Maximum structural simplex residual inherited from source & $8.15\times10^{-13}$ \\
    Maximum component discrepancy between independent evaluators & $3.31\times10^{-12}$ \\
    Difference between fourth- and eighth-power evaluations & $1.47\times10^{-12}$ \\
    Unexpected changed parameter groups & $0$ \\
    \bottomrule
  \end{tabular}
\end{table}

The safe certificate has SHA-256 digest
\begin{center}
\small\ttfamily
65d78064ef3686f850558b3881446403afe9d4910d2d2e046c3dab92c390ff83.
\end{center}
At its stored exponent, the two sides of the certificate inequality are
\[
  7.783640598318 \quad\text{and}\quad 4\log 7=7.783640596221,
\]
respectively.  The dual-safe reconstruction has logarithmic margin
$1.7464\times10^{-8}$ at the stored exponent.

\begin{proof}[Proof of \cref{thm:main}]
By \cref{lem:legal}, the 126 optimised row vectors are admissible parameters in the degeneration theorems
of~\cite{MoreAsymmetry}.  The fresh reconstruction evaluates all dependent complete-split and
compatibility quantities.  By \cref{lem:entropy-dual}, the dual entropy audit bounds every
maximum-entropy penalty in the conservative direction and yields
$\Omega(\rho)\le2.3713098963360215<2.371310$.  Substitution into
\cref{eq:master} and Sch\"onhage's asymptotic sum inequality therefore imply
$\omega<2.371310$.
\end{proof}

\begin{remark}[Floating-point status]\label{rem:floating}
The ancillary verification follows the numerical-certificate convention of the recent matrix
multiplication literature: it supplies a finite floating-point witness, independent evaluators, and a
headline rounded away from the computed value.  The entropy dual removes the only use of approximate
KKT optimality.  A fully formal interval-arithmetic replay would be a useful final layer; it is not claimed
here.  Until the certificate has also been reproduced externally, we regard the result as a world-record
candidate rather than an independently established published record.
\end{remark}

\section{The exact eighth-power lift}

The new row point is found in the fourth-power programme, but the requested eighth-power realisation is
exact and immediate.

\begin{proposition}[Tensor-square lift]\label{prop:square}
Suppose a fourth-power degeneration satisfies \cref{eq:master} for an exponent $\tau$.  Tensoring that
degeneration with itself gives an eighth-power degeneration satisfying the corresponding eighth-power
inequality for the same $\tau$.
\end{proposition}

\begin{proof}
For large $n$, suppose the fourth-power construction degenerates a direct sum of
$K_n=2^{o(n)}$ copies of $\CW_5^{\otimes4n}$ into a direct sum of $V_n$ copies of
$\langle A_n,B_n,C_n\rangle$.  Tensoring the entire degeneration with itself gives a degeneration of a
direct sum of $K_n^2=2^{o(n)}$ copies of $\CW_5^{\otimes8n}$ into $V_n^2$ copies of
$\langle A_n^2,B_n^2,C_n^2\rangle$.  Consequently the logarithmic retained-copy contribution and every
logarithmic matrix-size contribution double, while the asymptotic-rank target changes from $4\log7$ to
$8\log7$.  The sufficient inequality is therefore precisely the square of the fourth-power inequality.
\end{proof}

As an independent numerical check, we also evaluated the product construction in the 36 ordered pairs of
the six hashing orientations.  The result is
\[
  \omega^{(8)}_{\rm product}=2.3713098963357147,
\]
only $1.47\times10^{-12}$ above the fresh fourth-power reconstruction.  The maximum row and column
marginal residual is $3.63\times10^{-14}$.  An aggregate 36-cell evaluator gives
$2.3713098963357138$.  These computations are not needed for \cref{prop:square}, but they check the
eighth-power bookkeeping independently.

We additionally searched fixed-marginal orientation couplings and counterpart-specialised policies for the
80 highest-mass refined parents.  After rebasing on the final epigraph point, none improved the tensor
product.  Thus the present exponent improvement is a genuine fourth-power rebalancing inherited exactly
by the eighth power, not a claimed new correlated eighth-power construction.

\section{Discussion}

The improvement is numerically small but structurally informative.  The published programme already
contained the necessary freedom; what remained was a difficult nonsmooth balancing problem.  The
coordinate minima in \cref{eq:retained,eq:matrix} create narrow active-set intersections, so a smooth
surrogate or a generic large-scale SQP solve can stop near, but not at, the balance that controls the true
objective.  The explicit 42-inequality epigraph made those branches visible and allowed a targeted final
solve.

There are two immediate directions for further work.  First, one can repeat the active-set strategy while
releasing additional split distributions, rather than only the region weights.  Second, the full
$\CW_5^{\otimes8}$ programme permits non-product choices not represented by the exact tensor square.
Our exploratory searches found no gain on several natural restricted families, but they cover only a small
part of the full eighth-power feasible set.  A sparse reverse-mode implementation of the full programme,
combined with exact epigraph handling and certified entropy duals, remains a plausible route to another
improvement.

The broader lesson is methodological.  In record computations based on entropy and nested minima,
optimisation and certification should be separated.  A solver proposes structural parameters; a fresh
epigraph evaluation identifies the exact active branches; weak duality supplies conservative entropy
bounds; and an independent programme reconstructs the theorem's finite inequality.  This workflow is
cheap relative to the original parameter search and materially reduces the risk that a nominal improvement
is an artefact of smoothing, stale auxiliaries, or solver tolerances.

\section*{Acknowledgements and computational assistance}

The computational exploration, implementation of several independent verifiers, and preparation of the
initial manuscript used substantial assistance from OpenAI's Codex system based on GPT-5.6.  The system
was used to propose optimisation axes, test and reject unsuccessful eighth-power restrictions, derive the
entropy-dual audit, and cross-check the final certificate.  It is not listed as an author.  The mathematical
claims are presented in a form that can be checked independently from the ancillary files, and responsibility
for the submitted manuscript rests with the author.

\section{Reproducibility and ancillary files}

The principal files in the ancillary archive are listed in \cref{tab:ancillary}.

\begin{table}[t]
\centering
\caption{Principal machine-readable files supplied with the manuscript.}
\label{tab:ancillary}
\begin{tabular}{@{}>{\raggedright\arraybackslash}p{0.40\textwidth}>{\raggedright\arraybackslash}p{0.52\textwidth}@{}}
\toprule
File & Purpose \\
\midrule
\path{W1.00_2.371310_safe_final.mat} & Standalone 24,855-parameter safe certificate. \\
\path{p4_region_epigraph_stage4.npz} & The 126 optimised region rows and epigraph data. \\
\path{python_verifier.py} & Independent value-only reconstruction of the released model. \\
\path{dual_safe_verifier.py} & Entropy-dual conservative audit. \\
\path{p4_dense_epigraph_final_independent_verification.json} & Term-by-term dense audit. \\
\path{refined_aggregate_product_epigraph_final.json} & Explicit eighth-power product evaluation. \\
\path{SHA256SUMS} & Cryptographic digests for the archive. \\
\bottomrule
\end{tabular}
\end{table}

The authors' public archive contains the base implementation and published $2.371339$
certificate~\cite{MoreAsymmetryCode}.  The principal commands are listed below.
\begin{verbatim}
python3 python_verifier.py W1.00_2.371310_safe_final.mat --top 20
python3 dual_safe_verifier.py W1.00_2.371310_safe_final.mat
\end{verbatim}
The expected headline values are those in \cref{tab:values,tab:audit}.

\bibliographystyle{alpha}
\bibliography{references}

\end{document}
